\documentclass{beamer}
%\mode<presentation>
\usepackage{amsmath,amsfonts,amsthm}
\usepackage{tikz}

\newcommand{\E}{\mathbf{E}}
\newcommand{\F}{\mathbf{F}}
\newcommand{\Q}{\mathbb{Q}}
\newcommand{\R}{\mathbb{R}}
\newcommand{\Z}{\mathbb{Z}}

\newcommand{\cA}{\mathcal A}
\newcommand{\cB}{\mathcal B}
\newcommand{\cO}{\mathcal{O}}
\newcommand{\cH}{\mathcal H}
\newcommand{\C}{\mathbb C}
\newcommand{\cI}{\mathcal{I}}
\newcommand{\Rep}{\mathrm{Rep}}
\newcommand{\TVD}{\mathrm{TVD}}

\renewcommand{\>}{\rangle}
\newcommand{\<}{\langle}   

\begin{document}

%%%%%%%%%%%%%%%%%%%%%%%%%%%%%%%%%%%%%%%%%%%%%%%%%%%%%%%%%%%%%%%%%%%%%%%%%

\title{Quantum Algorithm for Computing the Period Lattice of an Infrastructure}
\author{Pawel M. Wocjan}
\institute{Department of Electrical Engineering and Computer Science\\University of Central Florida\\Orlando}
\date{}

%%%%%%%%%%%%%%%%%%%%%%%%%%%%%%%%%%%%%%%%%%%%%%%%%%%%%%%%%%%%%%%%%%%%%%%%%

\begin{frame}

\titlepage

\end{frame}

%%%%%%%%%%%%%%%%%%%%%%%%%%%%%%%%%%%%%%%%%%%%%%%%%%%%%%%%%%%%%%%%%%%%%%%%%

\begin{frame}{Joint work with Felix Fontein at the University of Z\"urich}

\begin{itemize}

\medskip
\item Based on 

\medskip
\begin{itemize}
\item \emph{Quantum Algorithms for Computing the Period Lattice of an Infrastructure}

 arXiv:1111.1348

\medskip
\item \emph{On the Probability of Generating a Lattice}

arXiv:1211.6246

\medskip
\item \emph{Computing Approximate Bases of a Lattice $L$ and Its Dual $L^*$ from an Approximate Generating Set of $L$}
\end{itemize}

\medskip
\item Supported by the NSF grant CCF-0726771 and the NSF CAREER Award CCF-0746600
\end{itemize}
\end{frame}


%%%%%%%%%%%%%%%%%%%%%%%%%%%%%%%%%%%%%%%%%%%%%%%%%%%%%%%%%%%%%%%%%%%%%%%%%

\begin{frame}{Quantum Algorithms and Hidden Lattice Problem over $\Z^n$}

\begin{itemize}
\item Efficient quantum algorithms provide the motivation for building a quantum computer

\medskip
\item Many quantum algorithms that achieve an exponential speed-up over classical algorithms can be formulated as a hidden lattice problem over $\Z^n$ 

\medskip
\item The hidden lattice problem over $\Z^n$ is very closely related to the hidden subgroup problem

\end{itemize}
\end{frame}

%%%%%%%%%%%%%%%%%%%%%%%%%%%%%%%%%%%%%%%%%%%%%%%%%%%%%%%%%%%%%%%%%%%%%%%%%

\begin{frame}{Definition of the Hidden Lattice Problem over $\Z^n$}

\begin{itemize}
\item
Let $A$ be a finite abelian group with generators $a_1,\ldots,a_n$

\medskip
\item Let $h : \Z^n \rightarrow A$ be a surjective group homomorphism given by
\[
h(x_1,\ldots,x_n) = a_1^{x_1} \cdots a_n^{x_n}
\] 

\item We assume that we can evaluate $h$ efficiently
\medskip
\item We have the isomorphism
\[
A \cong \Z^n / \Lambda \quad\mbox{where}\quad \Lambda=\ker h
\]

\item $\Lambda$ is a full-rank lattice in $\Z^n$

\medskip
\item The task is to determine a basis $\lambda_1,\ldots,\lambda_n$ of $\Z^n$
\[
\Lambda = \Z \lambda_1 + \ldots + \Z \lambda_n
\]
\end{itemize}

\end{frame}

%%%%%%%%%%%%%%%%%%%%%%%%%%%%%%%%%%%%%%%%%%%%%%%%%%%%%%%%%%%%%%%%%%%%%%%%%

\begin{frame}{Integer Factorization as Hidden Lattice Problem}

\begin{itemize}
\item Let $n$ be the integer that we seek to factor

\medskip
\item Choose an element $a$ uniformly at random from the unit group  
\[
\Z_n^\times=\{m \mid 1\le m < n, \; \gcd(m,n)=1\}
\]

\medskip
\item The surjective group homomorphism $h : \Z \rightarrow \Z_n^\times$ is given by 
\[
h(x) = a^x \pmod{n}
\]

\medskip
\item The kernel of $h$ is the lattice
\[
\Lambda = r \Z
\]
where $r$ is the order of $a$, i.e., the smallest positive integer $r$ such that $a^r=1 \pmod{n}$

\medskip
\item With high probability $r$ is even and $\gcd(a^{r/2}-1,n)$ yields a non-trivial factor
\end{itemize}

\end{frame}

%%%%%%%%%%%%%%%%%%%%%%%%%%%%%%%%%%%%%%%%%%%%%%%%%%%%%%%%%%%%%%%%%%%%%%%%%

\begin{frame}{Discrete Logarithm Problem as Hidden Lattice Problem}

\begin{itemize}
\item Let $a$ be a fixed generator of the unit group $\mathbb{F}_q^\times$ 

\medskip
\item Let $b$ a given element of $\mathbb{F}_q^\times$

\medskip
\item The task is to determine the discrete logarithm of $b$ with respect to the generator $a$, i.e., an integer $r\in\{0,\ldots,q-2\}$ such that 
\[
a^r=b
\]

\item The surjective homomorphism $h : \Z^2 \rightarrow \mathbb{F}_q^\times$ given by
\[
h(x,y)=b^x \cdot a^y
\]

\item The kernel of $f$ is the lattice $\Lambda$ generated by
\[
\Lambda = \Z \left(\begin{array}{c}0\\q-1\end{array}\right) + \Z \left(\begin{array}{c}1\\-r\end{array}\right)
\]
\end{itemize}
\end{frame}

%%%%%%%%%%%%%%%%%%%%%%%%%%%%%%%%%%%%%%%%%%%%%%%%%%%%%%%%%%%%%%%%%%%%%%%%%

%\begin{frame}{Some Algebraic Number Theory}

%\begin{itemize}
%\item Let $L$ be a number field of degree $d$ %, i.e., there exists a primitive element $\theta\in L$ such that $L=\Q[\theta]$ and the minimal polynomial of $\theta$ has degree $d$

%\item Let $\cO$ the ring of integers of $L$ %Let $\sigma_1,\ldots,\sigma_r$ the real embeddings of $L$ and $\rho_1,\ldots,\rho_

%\item Let $\sigma_1be the number of real embeddings and $2s$ be the number of complex embeddings of $L$ 

%\item The unit rank is $n=r+s-1$

%\item Dirichlet's Unit Theorem states that the unit group $\cO^\times$ of $L$ is isomorphic to
%\[
%R \times \Z^n 
%\]
%where $R$ is the group of roots of unity of $L$

%\item This proof relies on a logarithmic embedding of $L$ into $\R^n$
%\[
%\lambda(x) = \left( \sigma_1(x),\ldots,\sigma_r(x),\
%\]
%
%\end{itemize}
%\end{frame}

%%%%%%%%%%%%%%%%%%%%%%%%%%%%%%%%%%%%%%%%%%%%%%%%%%%%%%%%%%%%%%%%%%%%%%%%%

\begin{frame}{Quantum Algorithms and Hidden Lattice Problem over $\R^n$}

\begin{itemize}
\item The quantum algorithms due to {\sc Hallgren} and {\sc Vollmer and Schmidt}  for the number theoretic problems

\begin{itemize}
\item solving Pell's equation
\medskip
\item computing the unit group
\medskip
\item solving the principal ideal problem
\medskip
\item computing the class group 
\end{itemize}
can all be formulated as a hidden lattice problem over $\R^n$

\medskip
\item These running times of these quantum algorithms scales

\begin{itemize}
\item \textcolor{blue}{polynomially} in the logarithm of the discriminant $\Delta_{K/\Q}$ of the number field $K$, 

\medskip
\item but \textcolor{red}{exponentially} in $q(n)$ where $n$ is the unit rank of $K$ and $q : \mathbb{N} \rightarrow \mathbb{N}$ is a polynomial  
\end{itemize}
\end{itemize}
\end{frame}

%%%%%%%%%%%%%%%%%%%%%%%%%%%%%%%%%%%%%%%%%%%%%%%%%%%%%%%%%%%%%%%%%%%%%%%%%

\begin{frame}{Improvements}

\begin{itemize}
\item Achieve conceptual simplification by encapsulating all number theoretic details using infrastructures

\medskip
\item Provide a mathematically rigorous and complete analysis of the quantum algorithm

\medskip
\item Improve the running time by an exponential factor
\end{itemize}

\end{frame}

%%%%%%%%%%%%%%%%%%%%%%%%%%%%%%%%%%%%%%%%%%%%%%%%%%%%%%%%%%%%%%%%%%%%%%%%%

\begin{frame}{$n$-Dimensional Tori}

\begin{itemize}
\item Let $\Lambda$ be a full-rank lattice in $\R^n$, i.e.,

\[
\Lambda = \Z \lambda_1 + \ldots + \Z \lambda_n
\]

where $\lambda_1,\ldots,\lambda_n$ are linearly independent vectors in $\R^n$

\bigskip
\item The abelian group $T=\R^n/\Lambda$ is an $n$-dimensional torus

\bigskip
\item A one-dimensional torus $T=\R / (\Z \lambda)$ is simply a circle of circumference $\lambda$
\end{itemize}

\end{frame}

%%%%%%%%%%%%%%%%%%%%%%%%%%%%%%%%%%%%%%%%%%%%%%%%%%%%%%%%%%%%%%%%%%%%%%%%%

\begin{frame}{Hidden Lattice Problem over $\R^n$}

\begin{itemize}
\item Let $h : \R^n \rightarrow A$ be a surjective group homomorphism such that 
\[
\Lambda = \ker f
\] 
is a full-rank lattice in $\R^n$, i.e., $A$ is isomorphic to the $n$-dimensional torus 
\[
T=\R^n/\Lambda
\]

\medskip
\item Assume that $h$ can be evaluated ``approximately''

\medskip
\item The task is to find a $\gamma$-approximate basis of $\tilde{\lambda}_1,\ldots,\tilde{\lambda}_n$ of $\Lambda$

\[
\| \lambda_i - \tilde{\lambda}_i \|_2 \le \gamma \quad\mbox{for $i=1,\ldots,n$}
\]

\medskip
where $\lambda_1,\ldots,\lambda_n$ is some basis of $\Lambda$
\end{itemize}

\end{frame}


%%%%%%%%%%%%%%%%%%%%%%%%%%%%%%%%%%%%%%%%%%%%%%%%%%%%%%%%%%%%%%%%%%%%%%%%%

%\begin{frame}{One-dimensional Infrastructure}

%\begin{tikzpicture}
%\draw (0,0) [thick]circle (2);
%\foreach \x in {90, 80, 60, 30 , -30 , -110, -135, 150}
%\draw (0,0) +(\x:1.9) -- (\x:2.1);
%\draw (0,0) +(90:2) node[above] {$x_0$};
%\draw (0,0) +(84:2) node[above right] {$x_1$};
%\draw (0,0) +(-28:2) node[below right] {$x_i$};
%\end{tikzpicture}

%\end{frame}

%%%%%%%%%%%%%%%%%%%%%%%%%%%%%%%%%%%%%%%%%%%%%%%%%%%%%%%%%%%%%%%%%%%%%%%%%

\begin{frame}{Abstract Infrastructures -- Intuition}

\tikzstyle{background} = [dashed,white!67!black,fill=white!100!black]
  \tikzstyle{highlight} = [thick,fill=white!67!black]
  \tikzstyle{lines} = [black]
  \tikzstyle{fatlines} = [thick,black]
  \tikzstyle{dots} = [fill=black]
  \tikzstyle{fatdots} = [draw=black,thick,fill=black]
  
   \begin{center}
      \begin{tikzpicture}[y=-1cm, scale=0.45]
        % objects at depth 52:
        \draw[background] (9.5,1) rectangle (27.5,13.5);
       
        % objects at depth 50:
        \draw[lines] (15,10) -- (15,3.5) -- (19.5,3.5) -- (19.5,6) -- (21.5,6) -- (21.5,8.5) -- (22,8.5) -- (22,10) -- cycle;
        \draw[lines] (18.5,12.5) -- (23,12.5) -- (23,10.5) -- (22,10.5) -- (22,10) -- (18.5,10) -- cycle;
        \draw[lines] (22,10.5) rectangle (23,6);
        \draw[lines] (21.5,8.5) rectangle (22,6);
        \draw[lines] (19.5,6) -- (23,6) -- (23,5) -- (22.5,5) -- (22.5,3.5) -- (19.5,3.5) -- cycle;
        \draw[lines] (22.5,5) -- (24.5,5) -- (24.5,3.5) -- (23.5,3.5) -- (23.5,2) -- (22.5,2) -- cycle;
        \draw[lines] (11.5,7) -- (15,7) -- (15,5) -- (14.5,5) -- (14.5,4) -- (12.5,4) -- (12.5,3.5) -- (11.5,3.5) -- cycle;
        \draw[lines] (14.5,5) rectangle (15,3);
        \draw[lines] (12.5,4) rectangle (14.5,2.5);
        \draw[lines] (15,3.5) rectangle (17.5,3);
        \draw[lines] (17.5,3.5) -- (22.5,3.5) -- (22.5,2) -- (20.5,2) -- (20.5,1.5) -- (17.5,1.5) -- cycle;
        \draw[lines] (20.5,1) -- (20.5,2) -- (23.5,2) -- (23.5,1);
        \draw[lines] (16,1) -- (16,1.5) -- (17.5,1.5) -- (17.5,3) -- (14.5,3) -- (14.5,1);
        \draw[lines] (14.5,1) -- (14.5,2.5) -- (11.5,2.5) -- (11.5,1);
        \draw[lines] (10.5,1) -- (10.5,3.5) -- (11.5,3.5) -- (11.5,6) -- (9.5,6);
        \draw[lines] (25.5,1) -- (25.5,2) -- (26.5,2);
        \draw[lines] (26.5,1) -- (26.5,4.5) -- (27.5,4.5);
        \draw[lines] (24.5,3.5) -- (26.5,3.5) -- (26.5,4.5) -- (27.5,4.5);
        \draw[lines] (24.5,3.5) -- (24.5,7.5) -- (27.5,7.5);
        \draw[lines] (13.5,13.5) -- (13.5,7);
        \draw[lines] (20.5,12.5) -- (20.5,13.5);
        \draw[lines] (21,13.5) -- (21,12.5);
        \draw[lines] (23,12.5) -- (23,13) -- (27.5,13);
        
        \path[fatdots] (15,10) circle (0.1cm) node[below] {$x$};
        
        \path[fatdots] (19,7) circle (0.1cm);
        
        \draw[lines](15,10)--(19,7) node[below right] {\!\!\!\@$(x,t)$};
        
        %\draw (0,0) +(17,8) node[below right] {$(x,t)$};
        
        \draw[lines] (24.5,13.5) -- (24.5,13);
        \draw[lines] (26.5,13.5) -- (26.5,13);
      \end{tikzpicture}
   \end{center}
   \begin{itemize}
   \item The regions are ``cornered''

   \medskip
   \item The lower-left corners of the regions are labeled by elements $x$ of some finite set $X$
   
   \medskip
   \item Every point is uniquely identified by a pair $(x,t)$
   \end{itemize}
\end{frame}


%%%%%%%%%%%%%%%%%%%%%%%%%%%%%%%%%%%%%%%%%%%%%%%%%%%%%%%%%%%%%%%%%%%%%%%%%

\begin{frame}{Abstract Infrastructures -- Formal Definition}

An $n$-dimensional infrastructure $\mathcal{I}$ consists of

\begin{itemize}
\item a full-rank lattice $\Lambda$ in $\R^n$ (called the period lattice)

\medskip
\item a finite non-empty set $X$

\medskip
\item an injective map $d : X \rightarrow T$ where $T=\R^n/\Lambda$

\medskip
\item a set of $f$-representations $\mathrm{Rep}(\mathcal{I})\subseteq X\times \R^n$
such that 
\[
\Phi_\cI : \Rep(\cI) \rightarrow \R^n/\Lambda, \quad (x,t) \mapsto d(x) + (t \mbox{ mod } \Lambda)
\]
is a bijection

\end{itemize}
\end{frame}

%%%%%%%%%%%%%%%%%%%%%%%%%%%%%%%%%%%%%%%%%%%%%%%%%%%%%%%%%%%%%%%%%%%%%%%%%

\begin{frame}{Abstract Infrastructures -- Group Structure}

\begin{itemize}
\item $\Rep(\cI)$ is endowed with an abelian structure:
\[
(x,t) + (x',t') := \Phi_\cI^{-1}\big(\Phi_\cI(x,t) + \Phi_\cI(x',t')\big)
\]

\item There is a surjective group homomorphism $h$ given by
\[
h : \R^n \rightarrow \Rep(\cI), \quad v \mapsto \Phi_\cI^{-1}(v \mbox{ mod } \Lambda)
\]

\end{itemize}

\end{frame}

%%%%%%%%%%%%%%%%%%%%%%%%%%%%%%%%%%%%%%%%%%%%%%%%%%%%%%%%%%%%%%%%%%%%%%%%%

%\begin{frame}{Abstract Infrastructures}

%\begin{itemize}
%  \item An infrastructure~$\mathcal{I}$ gives rise to an abelian group isomorphic to the $n$-dimensional torus $T=\R^n/\Lambda$ 
%  where $\Lambda$ is a full-rank lattice in $\R^n$
  
%  \medskip
%  \item There is a finite set $X$ of distinguished points of $T$
  
%  \medskip
%  \item To each distinguished point $x$ of $T$, a region is assigned so that every point of $T$ lies in exactly one such region
  
%  \medskip
%  \item $\Rightarrow$ Every point~$y$ in the region of the distinguished point $x$ can be represented by the difference $t := y - x$ together
%  with $x$, i.e., as the pair $(x, t)$ 
  
%  \medskip
%  \item These tuples $(x, t)$ are called the $f$-representations of $\mathcal{I}$
%\end{itemize}
%
%\end{frame}

%%%%%%%%%%%%%%%%%%%%%%%%%%%%%%%%%%%%%%%%%%%%%%%%%%%%%%%%%%%%%%%%%%%%%%%%%

%\begin{frame}{Group Arithmetic with $f$-Representations}
%
%\begin{itemize}
%\item Denote the set of all $f$-representations by $\mathrm{Rep}(\mathcal{I})$
%
%\bigskip
%\item $\mathrm{Rep}(\mathcal{I})$ is a subset of $X \times \R^n$
%
%\bigskip
%\item $\mathrm{Rep}(\mathcal{I})$ is endowed with a group structure such that 
%\[
%\mathrm{Rep}(\mathcal{I}) \cong T = \R^n / \Lambda
%\]
%\end{itemize}
%
%\end{frame}

%%%%%%%%%%%%%%%%%%%%%%%%%%%%%%%%%%%%%%%%%%%%%%%%%%%%%%%%%%%%%%%%%%%%%%%%%

\begin{frame}{Concrete Infrastructures Arise From Number Fields}

\begin{itemize}
\item Let $K$ be a number field of degree $d=[K:\Q]$ and unit rank $n=r+s-1$ where $r$ is the number of real embeddings and $2s$ the number of complex embeddings of $K$
\medskip
\item The lattice $\Lambda$ is the kernel of the logarithmic map in Dirichlet's Unit theorem

\medskip
This theorem gives the structure of the unit group $\mathcal{O}_K^\times$ of the ring of integers $\mathcal{O}_K$ of $K$

\medskip
\item The set $X$ is the set of reduced fractional ideals of $\mathcal{O}_K$ 
\end{itemize}

\end{frame}

%%%%%%%%%%%%%%%%%%%%%%%%%%%%%%%%%%%%%%%%%%%%%%%%%%%%%%%%%%%%%%%%%%%%%%%%%

\begin{frame}{Perfect Periodicity is Impossible}

\begin{itemize}
\item The task is to determine an approximate basis of the period lattice $\Lambda$ hidden by the surjective group homomorphism

\[
h : \R^n \rightarrow \mathrm{Rep}(\mathcal{I})\subset X \times \R^n \quad \mbox{with}\quad
\ker f = \Lambda
\]

\item But there are two fundamental problems:

\medskip
\begin{itemize}
\item We can only evaluate $h$ at a finite number of points $v\in\R^n$

\medskip
\item We cannot compute \textcolor{blue}{the correct value $h(v)=(x,t)$}

\bigskip
We only obtain \textcolor{red}{an approximate $f$-representation $\tilde{h}(v)=(\tilde{x},\tilde{t})$} such that
\[
\Phi_\cI(x,t)  = \Phi_\cI(\tilde{x},\tilde{t}) + ( \delta \mbox{ mod } \Lambda)
\]

\medskip
where \textcolor{red}{the norm $\| \delta \|$} of the deviation $\delta\in\R^n$ is small

\end{itemize}

%\item $\Rightarrow$ We have to work with a finite number of carefully chosen evaluation points that avoid ``forbidden regions''
\end{itemize}

\end{frame}

%%%%%%%%%%%%%%%%%%%%%%%%%%%%%%%%%%%%%%%%%%%%%%%%%%%%%%%%%%%%%%%%%%%%%%%%%

\begin{frame}{Approximate $f$-representations}

\tikzstyle{background} = [dashed,white!67!black,fill=white!100!black]
  \tikzstyle{highlight} = [thick,fill=white!67!black]
  \tikzstyle{lines} = [black]
  \tikzstyle{fatlines} = [thick,black]
  \tikzstyle{dots} = [fill=black]
  \tikzstyle{fatdots} = [draw=black,thick,fill=black]
  \tikzstyle{approx} = [draw=red,thick]
  
   \begin{center}
      \begin{tikzpicture}[y=-1cm, scale=0.5]
        % objects at depth 52:
        \draw[background] (9.5,1) rectangle (27.5,13.5);
       
        % objects at depth 50:
        \draw[lines] (15,10) -- (15,3.5) -- (19.5,3.5) -- (19.5,6) -- (21.5,6) -- (21.5,8.5) -- (22,8.5) -- (22,10) -- cycle;
        \draw[lines] (18.5,12.5) -- (23,12.5) -- (23,10.5) -- (22,10.5) -- (22,10) -- (18.5,10) -- cycle;
        \draw[lines] (22,10.5) rectangle (23,6);
        \draw[lines] (21.5,8.5) rectangle (22,6);
        \draw[lines] (19.5,6) -- (23,6) -- (23,5) -- (22.5,5) -- (22.5,3.5) -- (19.5,3.5) -- cycle;
        \draw[lines] (22.5,5) -- (24.5,5) -- (24.5,3.5) -- (23.5,3.5) -- (23.5,2) -- (22.5,2) -- cycle;
        \draw[lines] (11.5,7) -- (15,7) -- (15,5) -- (14.5,5) -- (14.5,4) -- (12.5,4) -- (12.5,3.5) -- (11.5,3.5) -- cycle;
        \draw[lines] (14.5,5) rectangle (15,3);
        \draw[lines] (12.5,4) rectangle (14.5,2.5);
        \draw[lines] (15,3.5) rectangle (17.5,3);
        \draw[lines] (17.5,3.5) -- (22.5,3.5) -- (22.5,2) -- (20.5,2) -- (20.5,1.5) -- (17.5,1.5) -- cycle;
        \draw[lines] (20.5,1) -- (20.5,2) -- (23.5,2) -- (23.5,1);
        \draw[lines] (16,1) -- (16,1.5) -- (17.5,1.5) -- (17.5,3) -- (14.5,3) -- (14.5,1);
        \draw[lines] (14.5,1) -- (14.5,2.5) -- (11.5,2.5) -- (11.5,1);
        \draw[lines] (10.5,1) -- (10.5,3.5) -- (11.5,3.5) -- (11.5,6) -- (9.5,6);
        \draw[lines] (25.5,1) -- (25.5,2) -- (26.5,2);
        \draw[lines] (26.5,1) -- (26.5,4.5) -- (27.5,4.5);
        \draw[lines] (24.5,3.5) -- (26.5,3.5) -- (26.5,4.5) -- (27.5,4.5);
        \draw[lines] (24.5,3.5) -- (24.5,7.5) -- (27.5,7.5);
        \draw[lines] (13.5,13.5) -- (13.5,7);
        \draw[lines] (20.5,12.5) -- (20.5,13.5);
        \draw[lines] (21,13.5) -- (21,12.5);
        \draw[lines] (23,12.5) -- (23,13) -- (27.5,13);
        
        \path[fatdots] (15,10) circle (0.1cm) node[below] {$x$};
        
        \path[fatdots] (19,7) circle (0.1cm);
        
        \path[approx] (19,7) circle (2cm);
        
        \draw[lines](15,10)--(19,7) node[below right] {\!\!\!\@\textcolor{blue}{$(x,t)$}};
        
        \path[fatdots] (20,5.5) circle (0.1cm) node[above right] {\!\!\!\@\textcolor{red}{$(\tilde{x},\tilde{t})$}};
        
        %\draw (0,0) +(17,8) node[below right] {$(\tilde{x},\tilde{t})$};
        
        \draw[lines] (24.5,13.5) -- (24.5,13);
        \draw[lines] (26.5,13.5) -- (26.5,13);
      \end{tikzpicture}
   \end{center}
   \begin{itemize}
   \item The \textcolor{red}{circle} has radius $\|\delta\|$
   
   \end{itemize}
\end{frame}

%%%%%%%%%%%%%%%%%%%%%%%%%%%%%%%%%%%%%%%%%%%%%%%%%%%%%%%%%%%%%%%%%%%%%%%%%%%%

\begin{frame}{Pseudo-Periodicity is Possible}

\begin{itemize}
\item Let 
\begin{itemize}
\item $q$ and $N$ be sufficiently large positive integers

\item $\mathcal{V}=\{0,\ldots,qN-1\}^n$

\item $L$ an integer sufficiently greater than $N$

\item $s\in\{1,\ldots,L\}$ is chosen uniformly at random 
\end{itemize}

\medskip
\item Evaluate only at points of the form
\[
\frac{s}{L} + \frac{v}{N} \quad \mbox{for } v\in\mathcal{V}
\]

\item With constant probability, all evaluation points avoid certain \textcolor{orange}{forbidden regions} so that
\[
\big(x,\lfloor t / N\rfloor\big) = \big(\tilde{x},\lfloor \tilde{t} / N\rfloor\big)
\]
provided that \textcolor{red}{the approximation error $\|\delta\|$} is smaller than $\tfrac{1}{2}$ of the width of the forbidden regions

\medskip
$h(s/L+v/N)=(x,t)$ and $\tilde{h}(s/L+v/N)=(\tilde{x},\tilde{t})$
\end{itemize}

\end{frame}

%%%%%%%%%%%%%%%%%%%%%%%%%%%%%%%%%%%%%%%%%%%%%%%%%%%%%%%%%%%%%%%%%%%%%%%%%%%%

\begin{frame}{Forbidden regions I}

\tikzstyle{background} = [dashed,white!67!black,fill=white!100!black]
  \tikzstyle{highlight} = [thick,fill=white!67!orange]
  \tikzstyle{lines} = [black]
  \tikzstyle{fatlines} = [thick,black]
  \tikzstyle{dots} = [fill=black]
  \tikzstyle{fatdots} = [draw=black,thick,fill=black]
  
   \begin{center}
      \begin{tikzpicture}[y=-1cm, scale=0.5]
        % objects at depth 52:
        \draw[background] (9.5,1) rectangle (27.5,13.5);
        % objects at depth 51:
        \path[highlight] (10.4,1) rectangle (10.6,3.6);
        \path[highlight] (10.4,3.6) rectangle (12.6,3.4);
        \path[highlight] (12.4,4.1) rectangle (12.6,2.4);
        \path[highlight] (11.6,1) rectangle (11.4,2.6);
        \path[highlight] (11.4,2.6) rectangle (14.6,2.4);
        \path[highlight] (14.4,1) rectangle (14.6,5.1);
        \path[highlight] (12.4,4.1) rectangle (14.6,3.9);
        \path[highlight] (14.4,5.1) rectangle (15.1,4.9);
        \path[highlight] (9.5,6.1) rectangle (11.6,5.9);
        \path[highlight] (11.4,7.1) rectangle (11.6,3.4);
        \path[highlight] (11.4,7.1) rectangle (15.1,6.9);
        \path[highlight] (14.9,2.9) rectangle (15.1,10.1);
        \path[highlight] (14.4,2.9) rectangle (17.6,3.1);
        \path[highlight] (17.4,1.4) rectangle (17.6,3.6);
        \path[highlight] (15.9,1.6) rectangle (16.1,1);
        \path[highlight] (15.9,1.6) rectangle (20.6,1.4);
        \path[highlight] (20.4,1) rectangle (20.6,2.1);
        \path[highlight] (14.9,3.4) rectangle (22.6,3.6);
        \path[highlight] (20.4,2.1) rectangle (23.6,1.9);
        \path[highlight] (23.4,1) rectangle (23.6,3.6);
        \path[highlight] (22.4,1.9) rectangle (22.6,5.1);
        \path[highlight] (22.4,5.1) rectangle (24.6,4.9);
        \path[highlight] (19.4,6.1) rectangle (19.6,3.4);
        \path[highlight] (19.4,6.1) rectangle (23.1,5.9);
        \path[highlight] (23.4,3.6) rectangle (26.6,3.4);
        \path[highlight] (25.4,1) rectangle (25.6,2.1);
        \path[highlight] (25.4,2.1) rectangle (26.6,1.9);
        \path[highlight] (26.4,1) rectangle (26.6,4.6);
        \path[highlight] (26.4,4.6) rectangle (27.5,4.4);
        \path[highlight] (24.4,3.4) rectangle (24.6,7.6);
        \path[highlight] (24.4,7.6) rectangle (27.5,7.4);
        \path[highlight] (21.4,5.9) rectangle (21.6,8.6);
        \path[highlight] (21.4,8.6) rectangle (22.1,8.4);
        \path[highlight] (21.9,5.9) rectangle (22.1,10.6);
        \path[highlight] (21.9,10.6) rectangle (23.1,10.4);
        \path[highlight] (22.9,4.9) rectangle (23.1,13.1);
        \path[highlight] (22.9,13.1) rectangle (27.5,12.9);
        \path[highlight] (14.9,10.1) rectangle (22.1,9.9);
        \path[highlight] (18.4,9.9) rectangle (18.6,12.6);
        \path[highlight] (18.4,12.6) rectangle (23.1,12.4);
        \path[highlight] (13.4,6.9) rectangle (13.6,13.5);
        \path[highlight] (20.4,12.4) rectangle (20.6,13.5);
        \path[highlight] (20.9,13.5) rectangle (21.1,12.4);
        \path[highlight] (24.4,12.9) rectangle (24.6,13.5);
        \path[highlight] (26.4,13.5) rectangle (26.6,12.9);
        % objects at depth 50:
        \draw[lines] (15,10) -- (15,3.5) -- (19.5,3.5) -- (19.5,6) -- (21.5,6) -- (21.5,8.5) -- (22,8.5) -- (22,10) -- cycle;
        \draw[lines] (18.5,12.5) -- (23,12.5) -- (23,10.5) -- (22,10.5) -- (22,10) -- (18.5,10) -- cycle;
        \draw[lines] (22,10.5) rectangle (23,6);
        \draw[lines] (21.5,8.5) rectangle (22,6);
        \draw[lines] (19.5,6) -- (23,6) -- (23,5) -- (22.5,5) -- (22.5,3.5) -- (19.5,3.5) -- cycle;
        \draw[lines] (22.5,5) -- (24.5,5) -- (24.5,3.5) -- (23.5,3.5) -- (23.5,2) -- (22.5,2) -- cycle;
        \draw[lines] (11.5,7) -- (15,7) -- (15,5) -- (14.5,5) -- (14.5,4) -- (12.5,4) -- (12.5,3.5) -- (11.5,3.5) -- cycle;
        \draw[lines] (14.5,5) rectangle (15,3);
        \draw[lines] (12.5,4) rectangle (14.5,2.5);
        \draw[lines] (15,3.5) rectangle (17.5,3);
        \draw[lines] (17.5,3.5) -- (22.5,3.5) -- (22.5,2) -- (20.5,2) -- (20.5,1.5) -- (17.5,1.5) -- cycle;
        \draw[lines] (20.5,1) -- (20.5,2) -- (23.5,2) -- (23.5,1);
        \draw[lines] (16,1) -- (16,1.5) -- (17.5,1.5) -- (17.5,3) -- (14.5,3) -- (14.5,1);
        \draw[lines] (14.5,1) -- (14.5,2.5) -- (11.5,2.5) -- (11.5,1);
        \draw[lines] (10.5,1) -- (10.5,3.5) -- (11.5,3.5) -- (11.5,6) -- (9.5,6);
        \draw[lines] (25.5,1) -- (25.5,2) -- (26.5,2);
        \draw[lines] (26.5,1) -- (26.5,4.5) -- (27.5,4.5);
        \draw[lines] (24.5,3.5) -- (26.5,3.5) -- (26.5,4.5) -- (27.5,4.5);
        \draw[lines] (24.5,3.5) -- (24.5,7.5) -- (27.5,7.5);
        \draw[lines] (13.5,13.5) -- (13.5,7);
        \draw[lines] (20.5,12.5) -- (20.5,13.5);
        \draw[lines] (21,13.5) -- (21,12.5);
        \draw[lines] (23,12.5) -- (23,13) -- (27.5,13);
        \draw[lines] (24.5,13.5) -- (24.5,13);
        \draw[lines] (26.5,13.5) -- (26.5,13);
      \end{tikzpicture}
   \end{center}
   \begin{itemize}
   \item Regions I avoided $\Rightarrow x = \tilde{x}$
   \end{itemize}
\end{frame}

%%%%%%%%%%%%%%%%%%%%%%%%%%%%%%%%%%%%%%%%%%%%%%%%%%%%%%%%%%%%%%%%%%%%%%%%%%%%

\begin{frame}{Forbidden regions II}

\tikzstyle{background} = [dashed,white!67!black,fill=white!100!black]
  \tikzstyle{highlight} = [thick,fill=white!67!orange]
  \tikzstyle{lines} = [black]
  \tikzstyle{fatlines} = [thick,black]
  \tikzstyle{dots} = [fill=black]
  \tikzstyle{fatdots} = [draw=black,thick,fill=black]

    \begin{center}
      \begin{tikzpicture}[y=-1cm, scale=0.5]
        % objects at depth 53:
        \draw[background] (9.5,1) rectangle (27.5,13.5);
        % objects at depth 52:
        \path[highlight] (12.5,7) rectangle (12.7,4);
        \path[highlight] (13.6,7) rectangle (13.8,4);
        \path[highlight] (14.7,7) rectangle (14.9,5);
        \path[highlight] (11.5,6) rectangle (15,5.8);
        \path[highlight] (11.5,4.9) rectangle (14.5,4.7);
        \path[highlight] (11.5,3.8) rectangle (12.5,3.6);
        \path[highlight] (11.5,3.5) rectangle (11.7,2.5);
        \path[highlight] (10.5,2.5) rectangle (11.5,2.3);
        \path[highlight] (10.5,1.4) rectangle (11.5,1.2);
        \path[highlight] (13.5,4) rectangle (13.7,2.5);
        \path[highlight] (12.5,3) rectangle (14.4,2.8);
        \path[highlight] (14.5,4) rectangle (15,3.8);
        \path[highlight] (16,3.5) rectangle (16.2,3);
        \path[highlight] (17.1,3) rectangle (17.3,3.5);
        \path[highlight] (15.5,3) rectangle (15.7,1);
        \path[highlight] (16.6,3) rectangle (16.8,1.5);
        \path[highlight] (14.5,2) rectangle (17.5,1.8);
        \path[highlight] (17,1.5) rectangle (17.2,1);
        \path[highlight] (18.1,1.5) rectangle (18.3,1);
        \path[highlight] (19.2,1.4) rectangle (19.4,1);
        \path[highlight] (20.3,1.4) rectangle (20.5,1);
        \path[highlight] (18.5,3.5) rectangle (18.7,1.4);
        \path[highlight] (19.6,1.5) rectangle (19.8,3.5);
        \path[highlight] (20.7,3.4) rectangle (20.9,2);
        \path[highlight] (21.8,3.5) rectangle (22,2);
        \path[highlight] (17.5,2.5) rectangle (22.5,2.3);
        \path[highlight] (23.5,5) rectangle (23.7,3.5);
        \path[highlight] (22.5,4) rectangle (24.5,3.8);
        \path[highlight] (22.5,2.9) rectangle (23.5,2.7);
        \path[highlight] (21.5,2) rectangle (21.7,1);
        \path[highlight] (22.6,2) rectangle (22.8,1);
        \path[highlight] (24.5,3.5) rectangle (24.7,1);
        \path[highlight] (25.6,3.5) rectangle (25.8,2);
        \path[highlight] (23.5,2.5) rectangle (26.5,2.3);
        \path[highlight] (23.5,1.4) rectangle (25.5,1.2);
        \path[highlight] (26.5,3.5) rectangle (27.5,3.3);
        \path[highlight] (26.5,2.4) rectangle (27.5,2.2);
        \path[highlight] (26.5,1.3) rectangle (27.5,1.1);
        \path[highlight] (25.5,7.5) rectangle (25.7,3.5);
        \path[highlight] (26.6,7.5) rectangle (26.8,4.4);
        \path[highlight] (24.5,6.5) rectangle (27.5,6.3);
        \path[highlight] (24.5,5.4) rectangle (27.5,5.2);
        \path[highlight] (24.5,4.3) rectangle (26.5,4.1);
        \path[highlight] (20.5,6) rectangle (20.7,3.5);
        \path[highlight] (21.6,6) rectangle (21.8,3.5);
        \path[highlight] (22.7,6) rectangle (22.9,5);
        \path[highlight] (19.5,5) rectangle (22.5,4.8);
        \path[highlight] (19.5,3.9) rectangle (22.5,3.7);
        \path[highlight] (21.5,7.5) rectangle (22,7.3);
        \path[highlight] (21.5,6.4) rectangle (22,6.2);
        \path[highlight] (22,9.5) rectangle (23,9.3);
        \path[highlight] (22,8.4) rectangle (23,8.2);
        \path[highlight] (22,7.3) rectangle (23,7.1);
        \path[highlight] (23,6) rectangle (22,6.2);
        \path[highlight] (15,9) rectangle (22,8.8);
        \path[highlight] (15,7.9) rectangle (21.5,7.7);
        \path[highlight] (15,6.8) rectangle (21.5,6.6);
        \path[highlight] (15,5.5) rectangle (19.5,5.7);
        \path[highlight] (15,4.4) rectangle (19.5,4.6);
        \path[highlight] (16,10) rectangle (16.2,3.5);
        \path[highlight] (17.1,10) rectangle (17.3,3.5);
        \path[highlight] (18.2,9.9) rectangle (18.4,3.5);
        \path[highlight] (19.3,10) rectangle (19.5,3.5);
        \path[highlight] (20.4,10) rectangle (20.6,6);
        \path[highlight] (21.4,8.5) rectangle (21.6,10);
        \path[highlight] (19.5,12.5) rectangle (19.7,10);
        \path[highlight] (20.6,12.5) rectangle (20.8,10);
        \path[highlight] (21.7,12.5) rectangle (21.9,10);
        \path[highlight] (22.8,12.5) rectangle (23,10.5);
        \path[highlight] (18.5,11.5) rectangle (23,11.3);
        \path[highlight] (18.5,10.4) rectangle (22,10.2);
        \path[highlight] (24,13) rectangle (24.2,5);
        \path[highlight] (25.1,13) rectangle (25.3,7.5);
        \path[highlight] (26.2,13) rectangle (26.4,7.5);
        \path[highlight] (27.3,13) rectangle (27.5,7.5);
        \path[highlight] (23,12) rectangle (27.5,11.8);
        \path[highlight] (23,10.9) rectangle (27.5,10.7);
        \path[highlight] (23,9.8) rectangle (27.5,9.6);
        \path[highlight] (23,8.5) rectangle (27.5,8.7);
        \path[highlight] (23,7.6) rectangle (24.5,7.4);
        \path[highlight] (24.4,7.6) rectangle (27.5,7.5);
        \path[highlight] (23,6.5) rectangle (24.5,6.3);
        \path[highlight] (23,5.4) rectangle (24.5,5.2);
        \path[highlight] (25.5,13.5) rectangle (25.7,13);
        \path[highlight] (22,12.6) rectangle (22.2,13.5);
        \path[highlight] (23.1,13.5) rectangle (23.3,13);
        \path[highlight] (24.2,13.5) rectangle (24.4,13);
        \path[highlight] (21,13) rectangle (24.5,13.2);
        \path[highlight] (24.5,13.4) rectangle (26.5,13.5);
        \path[highlight] (20.5,13.1) rectangle (21,13.3);
        \path[highlight] (13.5,7.5) rectangle (15,7.7);
        \path[highlight] (15,8.8) rectangle (13.5,8.6);
        \path[highlight] (15,9.9) rectangle (13.5,9.7);
        \path[highlight] (13.5,11) rectangle (18.5,10.8);
        \path[highlight] (13.5,11.9) rectangle (18.5,12.1);
        \path[highlight] (13.5,13) rectangle (20.5,13.2);
        \path[highlight] (14.5,7) rectangle (14.7,13.5);
        \path[highlight] (15.6,10) rectangle (15.8,13.5);
        \path[highlight] (16.7,10) rectangle (16.9,13.5);
        \path[highlight] (18,10) rectangle (17.8,13.5);
        \path[highlight] (18.9,13.5) rectangle (19.1,12.5);
        \path[highlight] (20,13.5) rectangle (20.2,12.5);
        \path[highlight] (12.5,2.5) rectangle (12.7,1);
        \path[highlight] (13.6,2.5) rectangle (13.8,1);
        \path[highlight] (11.5,1.5) rectangle (14.5,1.3);
        \path[highlight] (10.8,6) rectangle (11,3.5);
        \path[highlight] (9.9,6) rectangle (9.7,1);
        \path[highlight] (11.5,5) rectangle (9.5,4.8);
        \path[highlight] (11.5,3.9) rectangle (9.5,3.7);
        \path[highlight] (9.5,2.6) rectangle (10.5,2.8);
        \path[highlight] (9.5,1.5) rectangle (10.5,1.7);
        \path[highlight] (12.9,7) rectangle (13.1,13.5);
        \path[highlight] (12,7) rectangle (11.8,13.5);
        \path[highlight] (10.9,6) rectangle (10.7,13.5);
        \path[highlight] (9.8,6) rectangle (9.6,13.5);
        \path[highlight] (11.5,6.7) rectangle (9.5,6.5);
        \path[highlight] (9.5,7.6) rectangle (13.5,7.8);
        \path[highlight] (9.5,8.9) rectangle (13.5,8.7);
        \path[highlight] (9.5,10) rectangle (13.5,9.8);
        \path[highlight] (13.6,10.9) rectangle (9.5,11.1);
        \path[highlight] (13.5,12) rectangle (9.5,12.2);
        \path[highlight] (13.5,13.1) rectangle (9.5,13.3);
        % objects at depth 51:
        \path[highlight] (10.4,1) rectangle (10.6,3.6);
        \path[highlight] (10.4,3.6) rectangle (12.6,3.4);
        \path[highlight] (12.4,4.1) rectangle (12.6,2.4);
        \path[highlight] (11.6,1) rectangle (11.4,2.6);
        \path[highlight] (11.4,2.6) rectangle (14.6,2.4);
        \path[highlight] (14.4,1) rectangle (14.6,5.1);
        \path[highlight] (12.4,4.1) rectangle (14.6,3.9);
        \path[highlight] (14.4,5.1) rectangle (15.1,4.9);
        \path[highlight] (9.5,6.1) rectangle (11.6,5.9);
        \path[highlight] (11.4,7.1) rectangle (11.6,3.4);
        \path[highlight] (11.4,7.1) rectangle (15.1,6.9);
        \path[highlight] (14.9,2.9) rectangle (15.1,10.1);
        \path[highlight] (14.4,2.9) rectangle (17.6,3.1);
        \path[highlight] (17.4,1.4) rectangle (17.6,3.6);
        \path[highlight] (15.9,1.6) rectangle (16.1,1);
        \path[highlight] (15.9,1.6) rectangle (20.6,1.4);
        \path[highlight] (20.4,1) rectangle (20.6,2.1);
        \path[highlight] (14.9,3.4) rectangle (22.6,3.6);
        \path[highlight] (20.4,2.1) rectangle (23.6,1.9);
        \path[highlight] (23.4,1) rectangle (23.6,3.6);
        \path[highlight] (22.4,1.9) rectangle (22.6,5.1);
        \path[highlight] (22.4,5.1) rectangle (24.6,4.9);
        \path[highlight] (19.4,6.1) rectangle (19.6,3.4);
        \path[highlight] (19.4,6.1) rectangle (23.1,5.9);
        \path[highlight] (23.4,3.6) rectangle (26.6,3.4);
        \path[highlight] (25.4,1) rectangle (25.6,2.1);
        \path[highlight] (25.4,2.1) rectangle (26.6,1.9);
        \path[highlight] (26.4,1) rectangle (26.6,4.6);
        \path[highlight] (26.4,4.6) rectangle (27.5,4.4);
        \path[highlight] (24.4,3.4) rectangle (24.6,7.6);
        \path[highlight] (24.4,7.6) rectangle (27.5,7.4);
        \path[highlight] (21.4,5.9) rectangle (21.6,8.6);
        \path[highlight] (21.4,8.6) rectangle (22.1,8.4);
        \path[highlight] (21.9,5.9) rectangle (22.1,10.6);
        \path[highlight] (21.9,10.6) rectangle (23.1,10.4);
        \path[highlight] (22.9,4.9) rectangle (23.1,13.1);
        \path[highlight] (22.9,13.1) rectangle (27.5,12.9);
        \path[highlight] (14.9,10.1) rectangle (22.1,9.9);
        \path[highlight] (18.4,9.9) rectangle (18.6,12.6);
        \path[highlight] (18.4,12.6) rectangle (23.1,12.4);
        \path[highlight] (13.4,6.9) rectangle (13.6,13.5);
        \path[highlight] (20.4,12.4) rectangle (20.6,13.5);
        \path[highlight] (20.9,13.5) rectangle (21.1,12.4);
        \path[highlight] (24.4,12.9) rectangle (24.6,13.5);
        \path[highlight] (26.4,13.5) rectangle (26.6,12.9);
        % objects at depth 50:
        \draw[lines] (15,10) -- (15,3.5) -- (19.5,3.5) -- (19.5,6) -- (21.5,6) -- (21.5,8.5) -- (22,8.5) -- (22,10) -- cycle;
        \draw[lines] (18.5,12.5) -- (23,12.5) -- (23,10.5) -- (22,10.5) -- (22,10) -- (18.5,10) -- cycle;
        \draw[lines] (22,10.5) rectangle (23,6);
        \draw[lines] (21.5,8.5) rectangle (22,6);
        \draw[lines] (19.5,6) -- (23,6) -- (23,5) -- (22.5,5) -- (22.5,3.5) -- (19.5,3.5) -- cycle;
        \draw[lines] (22.5,5) -- (24.5,5) -- (24.5,3.5) -- (23.5,3.5) -- (23.5,2) -- (22.5,2) -- cycle;
        \draw[lines] (11.5,7) -- (15,7) -- (15,5) -- (14.5,5) -- (14.5,4) -- (12.5,4) -- (12.5,3.5) -- (11.5,3.5) -- cycle;
        \draw[lines] (14.5,5) rectangle (15,3);
        \draw[lines] (12.5,4) rectangle (14.5,2.5);
        \draw[lines] (15,3.5) rectangle (17.5,3);
        \draw[lines] (17.5,3.5) -- (22.5,3.5) -- (22.5,2) -- (20.5,2) -- (20.5,1.5) -- (17.5,1.5) -- cycle;
        \draw[lines] (20.5,1) -- (20.5,2) -- (23.5,2) -- (23.5,1);
        \draw[lines] (16,1) -- (16,1.5) -- (17.5,1.5) -- (17.5,3) -- (14.5,3) -- (14.5,1);
        \draw[lines] (14.5,1) -- (14.5,2.5) -- (11.5,2.5) -- (11.5,1);
        \draw[lines] (10.5,1) -- (10.5,3.5) -- (11.5,3.5) -- (11.5,6) -- (9.5,6);
        \draw[lines] (25.5,1) -- (25.5,2) -- (26.5,2);
        \draw[lines] (26.5,1) -- (26.5,4.5) -- (27.5,4.5);
        \draw[lines] (24.5,3.5) -- (26.5,3.5) -- (26.5,4.5) -- (27.5,4.5);
        \draw[lines] (24.5,3.5) -- (24.5,7.5) -- (27.5,7.5);
        \draw[lines] (13.5,13.5) -- (13.5,7);
        \draw[lines] (20.5,12.5) -- (20.5,13.5);
        \draw[lines] (21,13.5) -- (21,12.5);
        \draw[lines] (23,12.5) -- (23,13) -- (27.5,13);
        \draw[lines] (24.5,13.5) -- (24.5,13);
        \draw[lines] (26.5,13.5) -- (26.5,13);
      \end{tikzpicture}
    \end{center}
    
    \begin{itemize}
   \item Regions II avoided $\Rightarrow$ not only $x = \tilde{x}$ but also $\lfloor t / N\rfloor=\lfloor \tilde{t} / N\rfloor$
   \end{itemize}

      
\end{frame}

%%%%%%%%%%%%%%%%%%%%%%%%%%%%%%%%%%%%%%%%%%%%%%%%%%%%%%%%%%%%%%%%%%%%%%%%%

\begin{frame}{Approximate Homomorphism $h_\approx$}

\begin{itemize}
\item Let $\mathcal{V}=\{0,\ldots,qN-1\}^n$ 

\medskip
\item We can efficiently compute a function $h_\approx : \mathcal{V} \rightarrow X\times\mathbb{N}$ with

\[
	h_\approx(v)=h_\approx(v') \quad\Longleftrightarrow\quad \frac{v-v'}{N} \approx \lambda \mbox{ for some $\lambda\in\Lambda$}
\]

\medskip
\item There is a collision iff $v$ and $v'$ differ approximately by an integer multiple of a lattice vector of the period lattice $\Lambda$

\medskip
\item The preimage $\mathcal{M}=h_{\approx}^{-1}(v)$ has the form
\[
	\mathcal{M} = \{ 
	v' = v + N \lambda + \xi_\lambda \in\mathcal{V} \mid \lambda\in\Lambda,\; \xi_\lambda\in(-1,1)^n
	\}
\]

provided that $s/L+v/N$ is not in a forbidden region III
\end{itemize}

\end{frame}


%%%%%%%%%%%%%%%%%%%%%%%%%%%%%%%%%%%%%%%%%%%%%%%%%%%%%%%%%%%%%%%%%%%%%%%%%

\begin{frame}{Pseudo-Lattice States}

\begin{itemize}
\item Let $\Lambda$ be a full-rank lattice in $\R^n$
and $\mathcal{V}=\{0,\ldots,qN-1\}^n$

\medskip
\item A uniform superposition in $\mathbb{C}^\mathcal{V}$ over a set of the form
\begin{align*}
\mathcal{M}  & = \{ 
v' = v + N \lambda + \xi_\lambda\in\mathcal{V} \mid \lambda\in\Lambda,\quad \xi_\lambda\in(-1,1)^n  
\}
\end{align*}

is called a pseudo-lattice state with period lattice $\Lambda$

\bigskip
\item We can efficiently prepare such pseudo-lattice states with constant probability by evaluating the approximate homomorphism $h_\approx$ 

\medskip
\item The offset $v$ is \textcolor{red}{random} and \textcolor{red}{different} each time we prepare a new pseudo-lattice state 

\medskip
\item The perturbations $\xi_\lambda$ spoil the lattice structure
 
\end{itemize}
\end{frame}

%%%%%%%%%%%%%%%%%%%%%%%%%%%%%%%%%%%%%%%%%%%%%%%%%%%%%%%%%%%%%%%%%%%%%%%%%

\begin{frame}{Quantum Algorithm}

\begin{itemize}

\item Prepare a pseudo-lattice state $|\mathcal{M}\>$ in $\mathbb{C}^{\mathcal{V}}$ with random offset $v$, i.e., a uniform superposition over a set of the form
\[
\mathcal{M}  = \{ 
v' = v + N \lambda + \xi_\lambda\in\mathcal{V} \mid \lambda\in\Lambda,\quad \xi_\lambda\in(-1,1)^n  \}
\]

\item Embed $|\mathcal{M}\>$ into the \textcolor{blue}{larger} Hilbert space $\mathbb{C}^\mathcal{W}$ where 
\[
\mathcal{W}=\{0,\ldots,\textcolor{blue}{2n}qN-1\}^n
\]
to deal with the perturbation effects caused by $\xi_\lambda$

\medskip
\item Apply the multidimensional QFT in $\C^\mathcal{W}$ to the``remove'' the random offset $v$ by using the shift-invariance property of the Fourier transform

\medskip
\item Measure to obtain outcome $w$

\medskip
\item The vector $w/(2nq)$ is an approximation of a lattice vector $\lambda^*$ of the dual lattice $\Lambda^*$ of $\Lambda$

\end{itemize}
\end{frame}

%%%%%%%%%%%%%%%%%%%%%%%%%%%%%%%%%%%%%%%%%%%%%%%%%%%%%%%%%%%%%%%%%%%%%%%%%
%%%%%%%%%%%%%%%%%%%%%%%%%%%%%%%%%%%%%%%%%%%%%%%%%%%%%%%%%%%%%%%%%%%%%%%%%
%%%%%%%%%%%%%%%%%%%%%%%%%%%%%%%%%%%%%%%%%%%%%%%%%%%%%%%%%%%%%%%%%%%%%%%%%

\begin{frame}{Sampling Approximate Vectors of the Dual Lattice $\Lambda^*$}

\begin{itemize}
\item Let $\Lambda^*$ denote the dual lattice of $\Lambda$, i.e.,

\[
\Lambda^* = \{ \lambda^*\in\R^n \mid \<\lambda^*,\lambda\>\in\Z \mbox{ for all } \lambda\in\Lambda \}
\]

\medskip
\item
We have to consider approximate vectors of $\Lambda^*$ since $\Lambda^*$ has non-integer coefficients

\medskip
\item For $\lambda^*=(\lambda^*_1,\ldots,\lambda^*_n)\in\Lambda^*$, define the set of good approximations by
\[
\mathcal{R}_{\lambda^*} = \big\{
(w_1,\ldots,w_n) \in\Z^n \mid w_k \in \{ \lfloor 2nq \lambda_k^* \rfloor, \lfloor 2nq \lambda_k^* \rfloor + 1 \} 
\big\}
\]

\medskip
\item For all $w\in\mathcal{R}_{\lambda^*}$, we have the approximation inequality

\[
\left\| \frac{w}{2 q n} - \lambda^* \right\|_2 \le \frac{1}{2\sqrt{n}q}
\]

%\medskip
%\quad each coordinate either rounded down or rounded up 

\end{itemize}
\end{frame}

%%%%%%%%%%%%%%%%%%%%%%%%%%%%%%%%%%%%%%%%%%%%%%%%%%%%%%%%%%%%%%%%%%%%%%%%%

\begin{frame}{Probability of Sampling an Approximate Vector of $\Lambda^*$}

\begin{itemize}
\item For all $\lambda^*\in\Lambda^*$ with $\mathcal{R}_{\lambda^*}\subset[0,2qn\kappa N]^n$ for some fixed $\kappa\in(0,1)$, we obtain the improved lower bound

\[
\Pr(\mathcal{R_{\lambda^*}} ) \gtrapprox  
2^{n-1} \cdot \frac{1}{(2qnN)^n} \cdot \frac{q^n}{\mathrm{det}(\Lambda)} \cdot \cos^2\big( \pi (\tfrac{1}{4} + \tfrac{1}{2qN} + 2\kappa n) \big) 
\]

\medskip
\item The number of such $\lambda^*$ $\approx$ $\tfrac{(\kappa N)^n}{\mathrm{det}(\Lambda^*)} = (\kappa N)^n \mathrm{det}(\Lambda)$

\medskip
\item The probability of obtaining an approximation of an essentially uniformly selected lattice vector $\lambda^* \in \Lambda^* \cap [0,\kappa N]^n$ is 

\[
p_1 \gtrapprox \frac{c}{2} \left(\frac{\kappa}{n}\right)^n
\]

\item The restriction to $[0,2qn\kappa N]^n$ with $\kappa<1$ allows us to deal with the perturbations effects caused by the deviations $(w/2qn) - \lambda^*$ for $w\in\mathcal{R}_{\lambda^*}$
\end{itemize}

\end{frame}

%%%%%%%%%%%%%%%%%%%%%%%%%%%%%%%%%%%%%%%%%%%%%%%%%%%%%%%%%%%%%%%%%%%%%%%%%

\begin{frame}{Probability of Generating a Full-Rank Lattice $L$}
\begin{itemize}
\item Let $L$ be a full-rank lattice in $\R^n$

\medskip
\item Assume $b_1,\ldots,b_n$ chosen uniformly at random in $L \cap [0,B)^n$ 

\medskip
\item With some constant probability these vectors generate a full-rank sublattice $L_0$

\medskip
\item $L/L_0$ is a finite abelian group that can be generated with only $n$ elements

\medskip
\item Assume $b_{n+1},\ldots,b_{2n+1}$ are chosen uniformly at random in $L \cap [0,B_0)^n$ with $B_0 > B$

\medskip
\item The elements $b_{n+k}+L_0$ are nearly uniformly distributioned over $L/L_0$ for $k=1,\ldots,n+1$

\medskip
\item With some constant probability $p$ the samples $b_1\ldots,b_{2n+1}$ generate $L$ 

\end{itemize}

\end{frame}

%%%%%%%%%%%%%%%%%%%%%%%%%%%%%%%%%%%%%%%%%%%%%%%%%%%%%%%%%%%%%%%%%%%%%%%%%

\begin{frame}{Approximate Basis for the Period Lattice $\Lambda$ (finally!)}

\begin{itemize}

\item Repeating the quantum algorithm $2n+1$ times yields the approximate generating set 

\[
	\frac{w_1}{2nq},\ldots,\frac{w_{2n+1}}{2nq} 
\]  

of $\Lambda^*$ with probability at least

\[
	p \, p^{2n+1} = p \left( \frac{c\kappa^n}{2 n^n} \right)^{2n+1} 
\]

where $c=\cos^2\big( \pi ( \tfrac{1}{4} + \tfrac{1}{2qN} + 2\kappa n) \big)$ and $\kappa \ge \tfrac{1}{8n}$

\medskip
\item We cover an approximate basis of $\Lambda$ from the above approximate generating set by a classical algorithm

\medskip
\item We recover an approximate basis of $\Lambda$ from the approximate basis of $\Lambda^*$ by matrix inversion
\end{itemize}


\end{frame}

%%%%%%%%%%%%%%%%%%%%%%%%%%%%%%%%%%%%%%%%%%%%%%%%%%%%%%%%%%%%%%%%%%%%%%%

\begin{frame}{Improvements}

\vspace{-0.2cm}
\begin{itemize}

\item We present a rigorous and complete proof that the running time is polynomial in the determinant of the period lattice $\Lambda$

\medskip
\item We give an explicit lower bound on the success probability 

\medskip
{\sc Hallgren} and {\sc Schmidt and Vollmer} did not give any lower bound

\medskip
{\sc Schmidt} presented a lower bound in his dissertation which is worse than our bound  

\medskip
\item We improve the success probability by at least $2^{n^2-1}$

\medskip
\item But the resulting lower bound is still quite bad $\Rightarrow$ find new ideas and techniques to improve it

\medskip

\end{itemize}

\end{frame}

\begin{frame}{Thank you!!!}

\end{frame}

\end{document} 


%%%%%%%%%%%%%%%%%%%%%%%%%%%%%%%%
%%%%%%%%%%%%%%%%%%%%%%%%%%%%%%%%
%%%%%%%%%%%%%%%%%%%%%%%%%%%%%%%%
%%%%%%%%%%%%%%%%%%%%%%%%%%%%%%%%
%%%%%%%%%%%%%%%%%%%%%%%%%%%%%%%%
%%%%%%%%%%%%%%%%%%%%%%%%%%%%%%%%
%%%%%%%%%%%%%%%%%%%%%%%%%%%%%%%%
%%%%%%%%%%%%%%%%%%%%%%%%%%%%%%%%
%%%%%%%%%%%%%%%%%%%%%%%%%%%%%%%%
%%%%%%%%%%%%%%%%%%%%%%%%%%%%%%%%
%%%%%%%%%%%%%%%%%%%%%%%%%%%%%%%%
%%%%%%%%%%%%%%%%%%%%%%%%%%%%%%%%
%%%%%%%%%%%%%%%%%%%%%%%%%%%%%%%%
%%%%%%%%%%%%%%%%%%%%%%%%%%%%%%%%
%%%%%%%%%%%%%%%%%%%%%%%%%%%%%%%%
%%%%%%%%%%%%%%%%%%%%%%%%%%%%%%%%
%%%%%%%%%%%%%%%%%%%%%%%%%%%%%%%%
%%%%%%%%%%%%%%%%%%%%%%%%%%%%%%%%
%%%%%%%%%%%%%%%%%%%%%%%%%%%%%%%%
%%%%%%%%%%%%%%%%%%%%%%%%%%%%%%%%
%%%%%%%%%%%%%%%%%%%%%%%%%%%%%%%%
%%%%%%%%%%%%%%%%%%%%%%%%%%%%%%%%
%%%%%%%%%%%%%%%%%%%%%%%%%%%%%%%%
%%%%%%%%%%%%%%%%%%%%%%%%%%%%%%%%
%%%%%%%%%%%%%%%%%%%%%%%%%%%%%%%%
%%%%%%%%%%%%%%%%%%%%%%%%%%%%%%%%
%%%%%%%%%%%%%%%%%%%%%%%%%%%%%%%%


